
\documentclass[12pt, a4paper]{article}

%-----------USEPACKAGE----------------------------
\usepackage{bm} % Tucne pismo
\usepackage[czech]{babel} % Cestina
\usepackage[T1]{fontenc}
\usepackage[utf8x]{inputenc}
\usepackage[unicode]{hyperref} % Odkazy v pdf, www a na e-mail
\usepackage{graphicx} % Obrazky
\usepackage{epstopdf} % Obrazky
\linespread{1.10} % Radkovani 1.3 odpovida radkovani 2
\usepackage{lmodern} % Daji se pouzit \HUGE atd.
\usepackage{amsmath}
\usepackage{algorithm}
\usepackage[noend]{algpseudocode} % Vkladani pseudo kodu
\usepackage[numbered, framed]{matlab-prettifier} % Vkladani kodu - vypada dobre
\lstset{style=Matlab-editor, basicstyle=\mlttfamily, escapechar=", mlshowsectionrules=true}
\graphicspath{{./img/}}




%------------------PROMENNE----------------------------
\newcommand{\cisloCviceni}{XX. cvičení}






%------------------LAYOUT----------------------------
\usepackage[top = 2.5 cm, bottom = 2.5 cm, left = 2.5 cm, right = 2.5 cm]{geometry} % geometrie stranky
\usepackage{longtable}% Pro dlouhy obsah, da se zalomit \pagebrek
\usepackage{fancyhdr}
\pagestyle{fancy}% Deffaultni nastaveni hlavicky a paticky
\setlength{\headheight}{16 pt}% Zvetsi hlavicku, aby to nedelalo warningy
\fancyhf{}
\lhead{Metody Počítačového Vidění}
\rhead{\cisloCviceni}
\fancyfoot[R]{\thepage}
\fancyfoot[L]{Verze 2.0, poslední úpravy: \today}





%--------------------------------------------------------
\begin{document}







 
%--------TITULNI-STRANA---------------------------------
\begin{titlepage}
\begin{center}
	\includegraphics[trim = 0.6cm 0.5cm 0.9cm 0.5cm, scale=1]{./FAV_logo_cz.pdf}
	\hspace*{\fill}
	\includegraphics[trim = 3.5cm 1.5cm 2.6cm 2cm, scale=0.295]{./KKY_logo_cz.pdf}\\
	\vspace*{\fill}
	\textbf{\Huge{Metody Počítačového Vidění \\ ~ \\ \cisloCviceni}}\\
	\vspace*{\fill}
	\textbf{\large{\href{mailto:LBures@kky.zcu.cz}{Ing. Lukáš Bureš}}} \hfill \textbf{\large{Plzeň, \today}}
\end{center}
\end{titlepage}








%--------OBSAH-CVICENI---------------------------------
\section*{Obsah cvičení}
\begin{itemize}
	\item 
	\item  
	\item 
	\item 
\end{itemize}









%----------------------------------------------------------
\section{Sekce}
\par{Sekce}

\subsection{Podsekce}
\par{Podsekce}









\end{document}



















